% Chapter 2: Background & related work — link prediction, structural/text baselines, DSAA 2023
\selectlanguage{greek}


\chapter{Θεωρητικό Υπόβαθρο και Σχετικές Εργασίες}
\label{ch:background}


\section{Γραφήματα με Κειμενικά Χαρακτηριστικά Κόμβων}
\label{sec:tag-graphs}

Ένα \textit{γράφημα με κειμενικά χαρακτηριστικά κόμβων} (\en{text-attributed graph,
TAG}) είναι ένα γράφημα $G = (V, E)$ στο οποίο κάθε κόμβος $v \in V$ συνδέεται με
ένα κειμενικό έγγραφο $d_v$.
Τέτοια γραφήματα απαντώνται ευρέως στην πράξη: τα
ακαδημαϊκά δίκτυα παραπομπών επισυνάπτουν περιλήψεις άρθρων στους κόμβους, τα
γραφήματα ηλεκτρονικού εμπορίου επισυνάπτουν περιγραφές προϊόντων, τα κοινωνικά δίκτυα
επισυνάπτουν προφίλ χρηστών, και εγκυκλοπαιδικά γραφήματα γνώσης όπως η \en{Wikipedia}
επισυνάπτουν το πλήρες κείμενο άρθρων.
Η πλούσια φύση των χαρακτηριστικών κόμβων
διαφοροποιεί τους \en{TAGs} από τα κοινά σχεσιακά γραφήματα και ανοίγει τον
δρόμο για τον συνδυασμό δομικής συλλογιστικής γραφημάτων με κατανόηση φυσικής γλώσσας.

Η πρόβλεψη ακμών σε \en{TAGs} απαιτεί μεθόδους που μπορούν να αξιοποιούν τόσο την
τοπολογία του $G$ όσο και τη σημασιολογία των $\{d_v\}$.
Οι καθαρά δομικές μέθοδοι
αγνοούν το κείμενο εξ ολοκλήρου και βασίζονται στην υπόθεση ότι η δομή της γειτονιάς
από μόνη της παρέχει επαρκές σήμα για την απόφαση ύπαρξης ακμής.
Η υπόθεση αυτή
παραβιάζεται συχνά σε αραιά ή εξελισσόμενα γραφήματα, όπου ένα μεγάλο κλάσμα
ζευγών κόμβων δεν μοιράζεται κοινούς γείτονες ακόμη και όταν ένας σύνδεσμος υπάρχει
ή θα έπρεπε να υπάρχει.
Αντίθετα, οι καθαρά κειμενικές μέθοδοι αγνοούν εντελώς το
γράφημα και ενδέχεται να συγχέουν εννοιολογικά συγγενείς κόμβους (π.χ., δύο άρθρα για
διαφορετικές μάχες του ίδιου πολέμου) με συνδεδεμένους.
Οι υβριδικές μέθοδοι που
συνδυάζουν αμφότερες τις τροπικότητες έχουν γίνει το κυρίαρχο παράδειγμα, αλλά
φέρουν το κόστος συμπερασμού τόσο ενός κωδικοποιητή γραφήματος όσο και ενός κωδικοποιητή
κειμένου, καθιστώντας δύσκολη την ανάπτυξή τους σε κλίμακα \en{Wikipedia}.

Το σύνολο αξιολόγησης \en{Wikipedia DSAA 2023}~\cite{phan2023link} αποτελεί ένα
κανονικό σύνολο δεδομένων πρόβλεψης ακμών σε \en{TAG}: το γράφημα έχει \GraphNodesTotal{} κόμβους
που αντιστοιχούν σε άρθρα της \en{Wikipedia}, και το κειμενικό έγγραφο κάθε κόμβου
είναι η περιγραφή της \en{Wikipedia} και προαιρετικά το περιεχόμενο του \en{infobox}.
Η παρούσα εργασία χρησιμοποιεί το \en{DSAA 2023} ως αρχικό πεδίο ανάπτυξης και
σύγκρισης με τον διαγωνισμό, και ένα ανεξάρτητα κατασκευασμένο σύνολο Απλής
Αγγλικής \en{Wikipedia} ως κύριο έλεγχο του \en{CascadeLP}.


\section{Δομικές Μέθοδοι}
\label{sec:structural-methods}

\subsection{Δομικοί Ευρετικοί Δείκτες}
\label{subsec:structural-heuristics}

Οι δομικοί ευρετικοί δείκτες ποσοτικοποιούν την τοπολογική εγγύτητα μεταξύ δύο κόμβων
χρησιμοποιώντας αποκλειστικά τη γειτνιακή δομή του γραφήματος.
Αποτελούν την αρχαιότερη
κατηγορία μεθόδων πρόβλεψης ακμών και παραμένουν ανταγωνιστικά μοντέλα αναφοράς
επειδή δεν απαιτούν δεδομένα εκπαίδευσης και υπολογίζονται τοπικά αφού
αποθηκευτούν οι λίστες γειτόνων.
Η \en{PA} απαιτεί μόνο ανάκτηση δύο βαθμών,
ενώ οι \en{CN}, \en{Jaccard} και \en{AA} απαιτούν τομή γειτονιών με κόστος που
εξαρτάται από τους βαθμούς των δύο άκρων.

Έστω $\Gamma(u)$ το σύνολο γειτόνων του κόμβου $u$ στο $G$.
Οι τέσσερις
ευρετικοί δείκτες που χρησιμοποιούνται στην παρούσα εργασία ορίζονται ως εξής.

\textbf{Κοινοί Γείτονες (\en{Common Neighbours, CN})}~\cite{newman2001clustering} μετράει τον αριθμό των κόμβων
που είναι γειτονικοί τόσο με τον $u$ όσο και με τον $v$:
\begin{equation}
    \mathrm{CN}(u, v) = |\Gamma(u) \cap \Gamma(v)|.
    \label{eq:cn}
\end{equation}
Η διαίσθηση είναι ότι δύο κόμβοι είναι πιθανότερο να συνδέονται εάν μοιράζονται ήδη
πολλές αμοιβαίες συνδέσεις, κατ' αναλογία με την κοινωνική αρχή ότι ένας φίλος
φίλου είναι πιθανότατα φίλος.

\textbf{Συντελεστής Jaccard (\en{Jaccard Coefficient})}~\cite{jaccard1912distribution} κανονικοποιεί τον αριθμό
κοινών γειτόνων με το μέγεθος της ένωσης των δύο γειτονιών, ελέγχοντας κόμβους
υψηλού βαθμού:
\begin{equation}
    \mathrm{Jaccard}(u, v) = \frac{|\Gamma(u) \cap \Gamma(v)|}{|\Gamma(u) \cup \Gamma(v)|}.
    \label{eq:jaccard}
\end{equation}
Ένα ζεύγος $(u, v)$ όπου $u$ είναι κόμβος υψηλού βαθμού με χιλιάδες γείτονες και $v$ έχει
μόνο δύο γείτονες θα λάβει χαμηλότερη βαθμολογία \en{Jaccard} από ένα ζεύγος κόμβων
χαμηλού βαθμού που μοιράζονται τον ίδιο αριθμό κοινών γειτόνων.

\textbf{Δείκτης \en{Adamic-Adar}}~\cite{adamic2003friends} σταθμίζει κάθε κοινό γείτονα
$w$ αντιστρόφως με τον λογάριθμο του βαθμού του, μειώνοντας τη βαρύτητα των κόμβων
υψηλού βαθμού που είναι κοινοί σχεδόν σε κάθε ζεύγος:
\begin{equation}
    \mathrm{AA}(u, v) = \sum_{w \in \Gamma(u) \cap \Gamma(v)} \frac{1}{\log |\Gamma(w)|}.
    \label{eq:aa}
\end{equation}
Ο δείκτης αυτός προτάθηκε αρχικά για τη μέτρηση της εγγύτητας δύο ατόμων στον
Παγκόσμιο Ιστό μετρώντας πόσες σελίδες συνδέονται και οι δύο, με έκπτωση για
σελίδες που συνδέονται από σχεδόν όλους.

\textbf{\en{Preferential Attachment} (PA)}~\cite{barabasi1999emergence}
παρακινείται από την παρατήρηση ότι σε πολλά πραγματικά δίκτυα νέες ακμές
συνδέονται κατά προτίμηση με κόμβους υψηλού βαθμού.
Η βαθμολογία είναι απλώς το
γινόμενο των βαθμών των δύο κόμβων:
\begin{equation}
    \mathrm{PA}(u, v) = |\Gamma(u)| \cdot |\Gamma(v)|.
    \label{eq:pa}
\end{equation}
Σε αντίθεση με τους άλλους τρεις δείκτες, το \en{PA} δεν απαιτεί υπολογισμό της τομής
συνόλων γειτόνων και είναι επομένως υπολογίσιμο σε $O(1)$ αφού προ-υπολογιστούν οι
βαθμοί.

Οι τέσσερις αυτοί δείκτες μοιράζονται έναν κοινό περιορισμό: εξετάζουν αποκλειστικά
την \textit{άμεση} γειτονιά δύο κόμβων (γείτονες σε απόσταση ενός βήματος) και
αγνοούν κάθε διαδρομή μεγαλύτερου μήκους μεταξύ τους.
Δύο παλαιότεροι δείκτες
αντιμετωπίζουν αυτόν τον περιορισμό εξετάζοντας διαδρομές αυθαίρετου μήκους.

Ο \textbf{Δείκτης \en{Katz}}~\cite{katz1953status} αθροίζει όλες τις διαδρομές μεταξύ
$u$ και $v$, σταθμίζοντας κάθε διαδρομή μήκους $l$ με έναν παράγοντα απόσβεσης
$\beta^l$ που μειώνει εκθετικά τη συμβολή μακρύτερων διαδρομών:
\begin{equation}
    \mathrm{Katz}(u, v) = \sum_{l=1}^{\infty} \beta^l \cdot |\mathrm{paths}_{u,v}^{\langle l \rangle}|
    = \left[(I - \beta A)^{-1} - I\right]_{uv},
    \label{eq:katz}
\end{equation}
όπου $A$ είναι ο πίνακας γειτνίασης του $G$, $I$ ο μοναδιαίος πίνακας, και
$|\mathrm{paths}_{u,v}^{\langle l \rangle}|$ το πλήθος διαδρομών μήκους $l$ μεταξύ
$u$ και $v$.
Η παράμετρος απόσβεσης $\beta$ πρέπει να είναι μικρότερη από το
αντίστροφο της μέγιστης ιδιοτιμής του $A$ για να συγκλίνει το άθροισμα.
Ο δείκτης \en{Katz} γενικεύει το \en{CN} (το οποίο αντιστοιχεί στον αριθμό
διαδρομών μήκους $l=2$)
με το κόστος ενός υπολογισμού αντιστροφής πίνακα, απαγορευτικού για γραφήματα
κλίμακας \en{Wikipedia} χωρίς προσεγγιστικές τεχνικές.

Το \textbf{\en{SimRank}}~\cite{jehwidom2002simrank} ορίζει την ομοιότητα δύο κόμβων
αναδρομικά, με τη διαίσθηση ότι «δύο αντικείμενα είναι όμοια εάν συνδέονται με
όμοια αντικείμενα»:
\begin{equation}
    \mathrm{sim}(u, v) = \frac{\gamma}{|\Gamma(u)| \, |\Gamma(v)|} \sum_{a \in \Gamma(u)} \sum_{b \in \Gamma(v)} \mathrm{sim}(a, b), \qquad u \neq v,
    \label{eq:simrank}
\end{equation}
με $\mathrm{sim}(u, u) = 1$ ως συνθήκη βάσης και $\gamma \in (0, 1)$ σταθερά απόσβεσης.
Η αναδρομή υπολογίζεται με επαναληπτική σταθεροποίηση (\en{fixed-point iteration})
πάνω σε όλα τα ζεύγη κόμβων ταυτόχρονα, κάτι που έχει πολυπλοκότητα $O(|V|^2)$ ανά
επανάληψη---επίσης απαγορευτικό στην πλήρη μορφή του για ένα γράφημα εκατοντάδων
χιλιάδων κόμβων, αλλά χρήσιμο ως εννοιολογικό προάγγελο των μεθόδων ενσωμάτωσης
γραφημάτων και ΝΔΓ που ακολουθούν, οι οποίες προσεγγίζουν την ίδια αναδρομική ιδέα
(η ομοιότητα ενός κόμβου εξαρτάται από την ομοιότητα των γειτόνων του) με
εκμαθημένες, υπολογιστικά τρακτάριστες αναπαραστάσεις.

Και οι έξι δείκτες καταρρέουν σε μηδέν (για \en{CN}, \en{Jaccard}, \en{AA}, \en{Katz} και \en{SimRank}) ή σε γινόμενο
βαθμών (για \en{PA}) όταν ένας ή και οι δύο κόμβοι δεν έχουν γείτονες στο γράφημα, γεγονός
που παρακινεί τον χειρισμό \en{cold-start} που περιγράφεται στο
Κεφάλαιο~\ref{ch:methodology}.

\subsection{Ενσωματώσεις Γραφημάτων}
\label{subsec:graph-embeddings}

Οι μέθοδοι ενσωμάτωσης γραφημάτων μαθαίνουν ένα χαμηλοδιάστατο συνεχές διάνυσμα
$\mathbf{z}_v \in \mathbb{R}^d$ για κάθε κόμβο $v$ έτσι ώστε ζεύγη κόμβων που είναι
δομικά ομοιόμορφοι στο $G$ να έχουν παρόμοιες ενσωματώσεις.\ Η πρόβλεψη ακμών
εκτελείται στη συνέχεια συνδυάζοντας τις ενσωματώσεις των δύο τελικών κόμβων,
τυπικά χρησιμοποιώντας το εσωτερικό γινόμενο, την απόσταση $L_1$ ή το γινόμενο
\en{Hadamard} ως χαρακτηριστικό ακμής.


Ενώ οι ενσωματώσεις γραφημάτων βελτιώνονται σε σχέση με τους ακατέργαστους ευρετικούς δείκτες
κωδικοποιώντας πληροφορία πολλαπλών βημάτων γειτονιάς, μοιράζονται τον θεμελιώδη
περιορισμό της αγνόησης του κειμένου κόμβων.
Απαιτούν επίσης επαναϋπολογισμό των
ενσωματώσεων όταν αλλάζει το γράφημα, κάτι που είναι δαπανηρό για τη \en{Wikipedia},
της οποίας η δομή υπερσυνδέσμων εξελίσσεται συνεχώς.

\subsection{Νευρωνικά Δίκτυα Γραφημάτων}
\label{subsec:gnns}

Οι μέθοδοι μη επαγωγικής ενσωμάτωσης γραφημάτων (π.χ. \en{DeepWalk}) μαθαίνουν μία σταθερή ενσωμάτωση ανά κόμβο,
αποθηκευμένη σε πίνακα αναζήτησης---δεν υπάρχει συνάρτηση που να χαρτογραφεί νέα
χαρακτηριστικά κόμβου σε ενσωμάτωση, οπότε η μέθοδος είναι εξ ορισμού \textit{μη
επαγωγική} (\en{transductive}): αδυνατεί να παράγει ενσωμάτωση για έναν κόμβο που δεν
είδε κατά την εκπαίδευση.
Τα \textbf{\en{Graph Neural Networks (GNNs)}} αντιμετωπίζουν
αυτόν τον περιορισμό μαθαίνοντας μια παραμετρική συνάρτηση \textit{μεταβίβασης
μηνυμάτων} (\en{message passing}) που μπορεί να εφαρμοστεί σε οποιονδήποτε κόμβο,
γνωστό ή νέο, εφόσον είναι γνωστή η γειτονιά του.

Το \textbf{\en{Graph Convolutional Network (GCN)}}~\cite{kipfwelling2017gcn} ορίζει
μια στρώση που ενημερώνει την αναπαράσταση κάθε κόμβου μέσω ενός σταθμισμένου μέσου
όρου των γειτόνων του:
\begin{equation}
    H^{(l+1)} = \sigma\!\left(\hat{D}^{-\frac{1}{2}} \hat{A} \hat{D}^{-\frac{1}{2}} H^{(l)} W^{(l)}\right),
    \label{eq:gcn}
\end{equation}
όπου $\hat{A} = A + I$ είναι ο πίνακας γειτνίασης με προστιθέμενους αυτοβρόχους ώστε
κάθε κόμβος να συμπεριλαμβάνει τη δική του προηγούμενη αναπαράσταση, $\hat{D}$ είναι ο
διαγώνιος πίνακας βαθμών του $\hat{A}$, $W^{(l)}$ είναι εκπαιδεύσιμος πίνακας βαρών
της στρώσης $l$, και $\sigma(\cdot)$ είναι μη γραμμική συνάρτηση ενεργοποίησης.
Η κανονικοποίηση $\hat{D}^{-1/2} \hat{A} \hat{D}^{-1/2}$ αποτρέπει κόμβους υψηλού
βαθμού να κυριαρχούν στο άθροισμα.
Στοιβάζοντας $L$ τέτοιες στρώσεις, κάθε κόμβος
ενσωματώνει πληροφορία από γείτονες έως $L$ βημάτων απόσταση.

Το \textbf{\en{GraphSAGE}}~\cite{hamilton2017graphsage} γενικεύει αυτή την ιδέα σε
ρητά \textit{επαγωγικό} (\en{inductive}) πλαίσιο: αντί να μαθαίνει σταθερό πίνακα
βαρών ανά κόμβο, μαθαίνει συναρτήσεις συγκέντρωσης γειτόνων που εφαρμόζονται
ομοιόμορφα σε ολόκληρο το γράφημα και άρα γενικεύονται σε κόμβους αόρατους κατά την
εκπαίδευση:
\begin{equation}
    h_v^{(l+1)} = \sigma\!\left(W^{(l)} \cdot \left[h_v^{(l)} \, \| \, \mathrm{AGG}^{(l)}\bigl(\{h_u^{(l)} : u \in \Gamma(v)\}\bigr)\right]\right),
    \label{eq:graphsage}
\end{equation}
όπου $\|$ δηλώνει συνένωση διανυσμάτων και $\mathrm{AGG}^{(l)}(\cdot)$ είναι μια
εκμαθημένη συνάρτηση συγκέντρωσης (μέσος όρος, \en{max-pooling}, ή \en{LSTM} πάνω σε
τυχαία διατεταγμένους γείτονες).
Αυτή η αρχιτεκτονική είναι αυτή που επιτρέπει σε
μεθόδους όπως τα \en{GraphFormers} (Ενότητα~\ref{sec:hybrid-architectures}) να
ενσωματώνουν πληροφορία γειτόνων εντός ενός Μετασχηματιστή κειμένου.

Το \textbf{\en{Graph Attention Network (GAT)}}~\cite{velickovic2018gat} αντικαθιστά
την ομοιόμορφη ή βαθμοεξαρτώμενη στάθμιση γειτόνων με εκμαθημένους συντελεστές
προσοχής, κατ' αναλογία με τον μηχανισμό \en{self-attention} που περιγράφεται
λεπτομερώς στην Ενότητα~\ref{sec:semantic-methods}:
\begin{equation}
    \alpha_{vu} = \frac{\exp\bigl(\mathrm{LeakyReLU}(\mathbf{a}^\top [W h_v \, \| \, W h_u])\bigr)}
    {\sum_{u' \in \Gamma(v)} \exp\bigl(\mathrm{LeakyReLU}(\mathbf{a}^\top [W h_v \, \| \, W h_{u'}])\bigr)},
    \label{eq:gat}
\end{equation}
όπου $\mathbf{a}$ και $W$ είναι εκπαιδεύσιμοι παράμετροι.
Έτσι, γείτονες που είναι
περισσότερο σχετικοί με τον κόμβο-στόχο συνεισφέρουν περισσότερο στην ενημέρωση της
αναπαράστασής του, αντί να μετρούν όλοι εξίσου.

Για την ίδια την εργασία της πρόβλεψης ακμών, το \textbf{\en{SEAL}}~\cite{zhangchen2018seal} συνδυάζει τη λογική των δεικτών \en{Katz}/\en{SimRank} με τα
\en{GNNs}: εξάγει ένα τοπικό υπο-γράφημα γύρω από το υποψήφιο ζεύγος $(u, v)$,
επισημαίνει κάθε κόμβο του υπο-γραφήματος με την απόσταση δικτύου του από $u$ και $v$
(\en{Double-Radius Node Labeling}), και εκπαιδεύει ένα \en{GNN} να ταξινομήσει
ολόκληρο το επισημασμένο υπο-γράφημα ως θετικό ή αρνητικό ζεύγος.
Αυτή η προσέγγιση
μαθαίνει αυτόματα ποια δομικά μοτίβα προβλέπουν καλά τις ακμές, αντί να επιλέγονται
εκ των προτέρων όπως στους ευρετικούς δείκτες, αλλά διατηρεί το κόστος εξαγωγής ενός
υπο-γραφήματος ανά υποψήφιο ζεύγος.
Όπως και τα \en{GCN}/\en{GraphSAGE}/\en{GAT}, το \en{SEAL} αγνοεί το κείμενο των
κόμβων εντελώς---αποτελεί καθαρά δομική μέθοδο, παρά τη χρήση νευρωνικού δικτύου.


\section{Σημασιολογικές Μέθοδοι}
\label{sec:semantic-methods}

\subsection{Κειμενικές Αναπαραστάσεις}
\label{subsec:text-representations}

Πριν από την εμφάνιση των προ-εκπαιδευμένων γλωσσικών μοντέλων, η πρόβλεψη ακμών
βασισμένη σε κείμενο βασιζόταν σε κλασικές αναπαραστάσεις επεξεργασίας φυσικής
γλώσσας.
Το \textbf{\en{TF-IDF}}~\cite{sparckjones1972statistical} (συχνότητα όρου-αντίστροφη συχνότητα εγγράφου)
αναπαριστά κάθε έγγραφο ως αραιό διάνυσμα όπου το βάρος του όρου $t$ στο έγγραφο
$d$ είναι ανάλογο της συχνότητάς του στο $d$ και αντιστρόφως ανάλογο του αριθμού
εγγράφων του σώματος που τον περιέχουν.
Ένα ζεύγος $(u, v)$ μπορεί να αναπαρασταθεί
από την ομοιότητα συνημίτονου των αντίστοιχων διανυσμάτων \en{TF-IDF}, από τη
συνένωσή τους ή από το στοιχειακό γινόμενό τους.
Το \en{TF-IDF} είναι υπολογιστικά
αποτελεσματικό αλλά αντιμετωπίζει το κείμενο ως μοντέλο \en{Bag-of-Words}, αγνοώντας τη σειρά
λέξεων και τις μορφολογικές σχέσεις.

\textbf{Τα χαρακτηριστικά συχνότητας μερών του λόγου} αναπαριστούν ένα έγγραφο
μετρώντας τη συχνότητα κάθε γραμματικής κατηγορίας (ουσιαστικό, ρήμα, επίθετο,
επίρρημα κ.λπ.) στο κείμενο.
Η διαίσθηση είναι ότι διαφορετικοί τύποι άρθρων
\en{Wikipedia} έχουν χαρακτηριστικές κατανομές μερών του λόγου: άρθρα για γεγονότα
είναι πλούσια σε ρήματα, άρθρα για τόπους είναι πλούσια σε ουσιαστικά με πολλά
κύρια ονόματα, και άρθρα για αφηρημένες έννοιες ενδέχεται να έχουν περισσότερα
επίθετα και προθέσεις.
Ένα 72-διάστατο διάνυσμα χαρακτηριστικών μερών του λόγου
που σχηματίζεται συνενώνοντας τα ιστογράμματα συχνότητας μερών του λόγου των δύο
τελικών κόμβων παρέχει ένα ελαφρύ αλλά εκπληκτικά πληροφοριακό σήμα για πρόβλεψη
ακμών σε αυτό το σύνολο αξιολόγησης, όπως κατέδειξαν οι Tran et al.~\cite{tran2023text}.

\subsection{Προ-εκπαιδευμένα Γλωσσικά Μοντέλα}
\label{subsec:pretrained-lms}

Ο πυρήνας κάθε μοντέλου \en{Transformer}---και άρα και των τριών προ-εκπαιδευμένων
μοντέλων που ακολουθούν---είναι ο μηχανισμός \en{self-attention} κλιμακούμενου
εσωτερικού γινομένου.
Δοθέντων πινάκων ερωτήματος (\en{query}) $Q$, κλειδιού
(\en{key}) $K$ και τιμής (\en{value}) $V$ που προκύπτουν από γραμμικές προβολές
των ενσωματώσεων εισόδου, η έξοδος προσοχής είναι:
\begin{equation}
    \mathrm{Attention}(Q, K, V) = \mathrm{softmax}\!\left(\frac{QK^\top}{\sqrt{d_k}}\right) V,
    \label{eq:self-attention}
\end{equation}
όπου $d_k$ είναι η διάσταση των διανυσμάτων κλειδιού, ο παράγοντας $1/\sqrt{d_k}$
αποτρέπει την κορεσμό της \en{softmax} για μεγάλες διαστάσεις, και το γινόμενο
$QK^\top$ παράγει μια βαθμολογία συνάφειας μεταξύ κάθε ζεύγους θέσεων \en{token}.
Πολλαπλές τέτοιες «κεφαλές προσοχής» (\en{multi-head attention}) εκτελούνται
παράλληλα σε διαφορετικούς υπόχωρους και συνενώνονται, επιτρέποντας στο μοντέλο
να καταγράφει ταυτόχρονα πολλαπλούς τύπους σχέσεων μεταξύ \en{tokens}.

\textbf{\en{BERT}}~\cite{devlin2019bert} (\en{Bidirectional Encoder Representations from
Transformers}) εισήγαγε το παράδειγμα της προ-εκπαίδευσης ενός βαθύ Μετασχηματιστή
σε μεγάλα κειμενικά σώματα χρησιμοποιώντας \en{Masked Language Modeling} (\en{MLM})
και πρόβλεψη επόμενης πρότασης (\en{NSP}). Στο \en{MLM}, ένα τυχαίο 15\% των \en{tokens}
στην είσοδο αντικαθίσταται με ένα \en{token} \texttt{\en{[MASK]}}, και το μοντέλο
εκπαιδεύεται να ελαχιστοποιεί την αρνητική λογαριθμική πιθανοφάνεια πρόβλεψης του
αρχικού \en{token} από το διπλής κατεύθυνσης πλαίσιο:
\begin{equation}
    \mathcal{L}_{\mathrm{MLM}} = -\frac{1}{|M|} \sum_{i \in M} \log P(x_i \mid x_{\setminus M}; \theta),
    \label{eq:bert-mlm}
\end{equation}
όπου $M$ είναι το σύνολο των καλυμμένων θέσεων και $x_{\setminus M}$ είναι η
ακολουθία εισόδου με τις θέσεις αυτές καλυμμένες.
Στο \en{NSP}, το μοντέλο εκπαιδεύεται με απώλεια δυαδικής διασταυρούμενης εντροπίας
να προβλέπει εάν δύο προτάσεις είναι συνεχόμενες στο αρχικό σώμα, χρησιμοποιώντας
την ενσωμάτωση του ειδικού \en{token} \texttt{\en{[CLS]}} ως είσοδο σε έναν
γραμμικό ταξινομητή.
Αυτή η προ-εκπαίδευση εμπεδώνει πλούσια συντακτική και σημασιολογική
γνώση που μπορεί να μεταφερθεί σε εργασίες τελικού σταδίου με προσαρμογή σε επισημαντικά
δεδομένα.
Ο αμφίδρομος μηχανισμός προσοχής της Εξίσωσης~\ref{eq:self-attention}, που επιτρέπει σε κάθε \en{token}
να προσέχει όλα τα άλλα \en{tokens} και στις δύο κατευθύνσεις, αντιτίθεται στα
μονόδρομα μοντέλα όπως το \en{GPT} και δίνει πλεονέκτημα στο \en{BERT} σε εργασίες που
απαιτούν κατανόηση του πλήρους πλαισίου.

\textbf{\en{RoBERTa}}~\cite{liu2019roberta} (\en{Robustly Optimised BERT Pre-training
Approach}) κατάργησε τον στόχο \en{NSP}, χρησιμοποίησε δυναμική κάλυψη (ώστε διαφορετικά
\en{tokens} να καλύπτονται σε κάθε διέλευση εκπαίδευσης), εκπαιδεύτηκε σε περισσότερα
δεδομένα με μεγαλύτερα μεγέθη δέσμης, και διαπίστωσε ότι αυτές οι αλλαγές βελτίωναν
συστηματικά την απόδοση σε εργασίες τελικού σταδίου.
Το \en{RoBERTa} έγινε το τυπικό ισχυρό
μοντέλο αναφοράς για εργασίες \en{NLP} επιπέδου πρότασης.

\textbf{\en{Sentence-BERT}}~\cite{reimers2019sentence} αντιμετώπισε τον περιορισμό ότι το
κανονικό \en{BERT} είναι αργό για εργασίες σημασιολογικής ομοιότητας επειδή απαιτεί έναν
\en{cross-encoder} που επεξεργάζεται αμφότερες τις προτάσεις από κοινού.
Το
\en{Sentence-BERT} εισάγει μια αρχιτεκτονική \en{Siamese} δικτύου στην οποία το ίδιο
μοντέλο \en{BERT} κωδικοποιεί κάθε πρόταση ανεξάρτητα σε ενσωμάτωση σταθερού μεγέθους,
και η ομοιότητα υπολογίζεται ως η απόσταση συνημίτονου μεταξύ των δύο ενσωματώσεων.
Αυτός ο διπλός κωδικοποιητής μειώνει τον χρόνο συμπερασμού από $O(n^2)$
(\en{cross-encoder} σε όλα τα ζεύγη) σε $O(n)$ (κωδικοποίηση κάθε κόμβου
μία φορά) συν $O(n^2)$ υπολογισμούς ομοιότητας συνημίτονου, κάτι που είναι πολύ
φθηνότερο όταν οι ενσωματώσεις μπορούν να προ-υπολογιστούν.
Το \en{Sentence-BERT}
προσαρμόζεται σε σύνολα δεδομένων \en{Natural Language Inference (NLI)} με απώλεια
\en{softmax} πάνω στη συνένωση $[\mathbf{u}; \mathbf{v}; |\mathbf{u} - \mathbf{v}|]$
των δύο ενσωματώσεων πρότασης $\mathbf{u}, \mathbf{v}$:
\begin{equation}
    o = \mathrm{softmax}\!\left(W_t \, [\mathbf{u}; \mathbf{v}; |\mathbf{u} - \mathbf{v}|]\right),
    \label{eq:sbert-softmax}
\end{equation}
όπου $W_t$ είναι εκπαιδεύσιμος πίνακας βαρών και οι τρεις κλάσεις εξόδου αντιστοιχούν
σε συνεπαγωγή, ουδετερότητα και αντίφαση· εναλλακτικά, εκπαιδεύεται με απώλεια
τύπου \en{contrastive/triplet} που ελαχιστοποιεί απευθείας την απόσταση
$\| \mathbf{u} - \mathbf{v} \|$ για σημασιολογικά όμοια ζεύγη και τη μεγιστοποιεί
(έως ένα περιθώριο $\epsilon$) για ανόμοια ζεύγη, φέρνοντας σημασιολογικά παρόμοια
προτάσεις πλησιέστερα στον χώρο ενσωμάτωσης χωρίς να απαιτείται ετικέτα κλάσης
\en{NLI}.


\section{Υβριδικές Αρχιτεκτονικές}
\label{sec:hybrid-architectures}

Οι υβριδικές μέθοδοι συνδυάζουν τοπολογία γραφήματος με κειμενικό περιεχόμενο, στοχεύοντας
στην ταυτόχρονη αξιοποίηση αμφότερων των πηγών σήματος.
Αρκετές αρχιτεκτονικές έχουν
προταθεί ειδικά για γραφήματα με κειμενικά χαρακτηριστικά κόμβων.

\textbf{\en{LinkBERT}}~\cite{yasunaga2022linkbert} επαυξάνει την προ-εκπαίδευση του \en{BERT}
αντικαθιστώντας το \en{NSP} με εργασία \textit{πρόβλεψης σχέσης εγγράφων} (\en{DRP}): αντί να
προβλέπει εάν δύο τυχαία επιλεγμένες προτάσεις είναι συνεχόμενες, το μοντέλο
εκπαιδεύεται σε τριτασικη ταξινόμηση της σχέσης $r$ μεταξύ δύο τμημάτων κειμένου
$d_1, d_2$:
\begin{equation}
    \mathcal{L}_{\mathrm{DRP}} = -\log P(r \mid d_1, d_2; \theta), \quad
    r \in \{\text{συνεχόμενο}, \text{τυχαίο}, \text{συνδεδεμένο}\},
    \label{eq:linkbert-drp}
\end{equation}
όπου η κλάση «συνδεδεμένο» ισχύει όταν τα $d_1$ και $d_2$ συνδέονται με
υπερσύνδεσμο στο γράφημα πηγής, «συνεχόμενο» όταν προέρχονται από το ίδιο έγγραφο,
και «τυχαίο» όταν προέρχονται από άσχετα έγγραφα.
Το σώμα
προ-εκπαίδευσης κατασκευάζεται δειγματοληπτώντας τέτοια ζεύγη εγγράφων από το γράφημα
υπερσυνδέσμων της \en{Wikipedia}.
Αυτό επιτρέπει στο μοντέλο να μαθαίνει
διαεγγραφικές σχέσεις που είναι αόρατες στην τυπική εκπαίδευση \en{BERT}, με ελάχιστο
πρόσθετο υπολογιστικό κόστος έναντι του \en{NSP} αφού η αρχιτεκτονική Μετασχηματιστή
παραμένει αμετάβλητη---μόνο η εργασία προ-εκπαίδευσης αλλάζει.
Το \en{LinkBERT}
ανέδειξε βελτίωση της απόδοσης σε εργασίες απάντησης ερωτήσεων έντασης γνώσης και
πολυβηματικής συλλογιστικής.

\textbf{\en{GraphFormers}}~\cite{yang2021graphformers} ενσωματώνουν έναν κωδικοποιητή
κειμένου Μετασχηματιστή εντός βρόχου μεταβίβασης μηνυμάτων ΝΔΓ, αντί να
επεξεργάζονται πρώτα το κείμενο και έπειτα το γράφημα ως δύο διαδοχικά, ανεξάρτητα
στάδια.
Σε κάθε επίπεδο $l$, η αναπαράσταση \en{token} του κόμβου $v$ ενημερώνεται
όχι μόνο από το εσωτερικό \en{self-attention} του δικού του κειμένου αλλά και
από μια συγκεντρωτική αναπαράσταση γειτόνων που εισάγεται ως πρόσθετο
\en{token} \en{[GNN]} στην ακολουθία εισόδου:
\begin{equation}
    h_v^{(l+1)} = \mathrm{Transformer}^{(l)}\!\left(h_v^{(l)}, \; \mathrm{AGG}\bigl(\{h_u^{(l)} : u \in \Gamma(v)\}\bigr)\right),
    \label{eq:graphformers-update}
\end{equation}
όπου $\mathrm{AGG}(\cdot)$ είναι μια συνάρτηση συγκέντρωσης της ίδιας οικογένειας με
αυτές του \en{GraphSAGE} (Εξίσωση~\ref{eq:graphsage})---μέσος όρος, \en{max-pooling},
ή στάθμιση τύπου \en{GAT} (Εξίσωση~\ref{eq:gat})---πάνω στις αναπαραστάσεις των
γειτόνων $\Gamma(v)$ του προηγούμενου επιπέδου.
Επαναλαμβάνοντας αυτή τη διαδικασία σε πολλαπλά επίπεδα, η κωδικοποίηση
κειμένου κάθε κόμβου ενσωματώνει προοδευτικά πληροφορία από γείτονες πολλαπλών
βημάτων, επιτρέποντας τη συγχώνευση δομικής και σημασιολογικής πληροφορίας σε
πολλαπλά επίπεδα αφαίρεσης αντί σε ένα μόνο τελικό στάδιο συνδυασμού.
Αυτή η αρχιτεκτονική επιτυγχάνει ισχυρή απόδοση σε εργασίες ταξινόμησης κόμβων
και πρόβλεψης ακμών σε ακαδημαϊκά δίκτυα παραπομπών, αλλά απαιτεί πρόσβαση στο
πλήρες γράφημα κατά τον χρόνο συμπερασμού (για τον υπολογισμό $\Gamma(v)$ σε κάθε
επίπεδο) και είναι δαπανηρή στην εκπαίδευση, αφού κάθε βήμα οπισθοδιάδοσης
διαδίδεται μέσω τόσο του Μετασχηματιστή όσο και του γραφήματος γειτόνων.

\textbf{\en{Patton}}~\cite{jin2023patton} προτείνει δύο στρατηγικές προ-εκπαίδευσης για
γραφήματα με κειμενικά χαρακτηριστικά κόμβων.
Η πρώτη, \en{Masked Language Modeling}
πλαισιοποιημένη από δίκτυο, επεκτείνει την Εξίσωση~\ref{eq:bert-mlm} ώστε η
πρόβλεψη ενός καλυμμένου \en{token} στο έγγραφο $d_v$ να εξαρτάται και από το
κείμενο των γειτόνων $\Gamma(v)$:
\begin{equation}
    \mathcal{L}_{\mathrm{MLM}}^{\mathrm{graph}} = -\frac{1}{|M|} \sum_{i \in M} \log P\bigl(x_i \mid x_{\setminus M}, \{d_u : u \in \Gamma(v)\}; \theta\bigr).
    \label{eq:patton-mlm}
\end{equation}
Η δεύτερη, \en{πρόβλεψη καλυμμένου κόμβου} (\en{Masked Node Prediction}), είναι μια
δυαδική εργασία διάκρισης που εκπαιδεύει το μοντέλο να προβλέπει εάν ένας υποψήφιος
κόμβος $u$ ανήκει πράγματι στη γειτονιά $\Gamma(v)$ ενός κόμβου-αγκύρωσης $v$,
βασιζόμενο αποκλειστικά στο κείμενό του $d_u$ (χωρίς άμεση πρόσβαση στην τοπολογία
κατά τον χρόνο συμπερασμού)---ανάλογη διάκριση θετικού/αρνητικού ζεύγους με αυτήν
που αντιμετωπίζει η παρούσα εργασία στο σύνολο \en{DSAA 2023}.
Το \en{Patton} παρέχει έτσι τόσο δομικό όσο και σημασιολογικό σήμα
προ-εκπαίδευσης, και αξιολογήθηκε σε ταξινόμηση άρθρων και πρόβλεψη ακμών στο
δίκτυο παραπομπών \en{OGB}.

Ο Πίνακας~\ref{tab:method-family-comparison} συνοψίζει τις οικογένειες μεθόδων που
παρουσιάστηκαν σε αυτό το κεφάλαιο, ως προς το εάν αξιοποιούν δομή γραφήματος ή/και
κείμενο, το εάν απαιτούν εκπαίδευση, και το κόστος συμπερασμού ανά ζεύγος---η
βασική διάσταση γύρω από την οποία οργανώνεται το \en{CascadeLP} στο
Κεφάλαιο~\ref{ch:methodology}.

\begin{table}[htbp]
    \centering
    \small
    \begin{tabular}{@{}p{3.1cm}ccp{2.4cm}p{4.2cm}@{}}
        \toprule
        Μέθοδος & Γράφημα & Κείμενο & Απαιτεί εκπαίδευση & Κόστος συμπερασμού ανά ζεύγος \\
        \midrule
        \en{CN} / \en{Jaccard} / \en{AA} / \en{PA} & \checkmark & -- & Όχι & $O(1)$--$O(|\Gamma|)$ \\
        \en{Katz} / \en{SimRank} & \checkmark & -- & Όχι & $O(|V|^3)$--$O(|V|^2)$ (πλήρης) \\
        \en{GCN} / \en{GraphSAGE} / \en{GAT} & \checkmark & -- & Ναι & $O(d)$ ανά επίπεδο, απαιτεί γειτονιά \\
        \en{SEAL} & \checkmark & -- & Ναι & Εξαγωγή $+$ ταξινόμηση υπο-γραφήματος ανά ζεύγος \\
        \en{TF-IDF} + ταξινομητής & -- & \checkmark & Ναι & $O(|V_{\text{λεξιλόγιο}}|)$ \\
        \en{POS} + Τυχαίο Δάσος & -- & \checkmark & Ναι & Γραμμικό στο μήκος κειμένου \\
        \en{BERT}/\en{RoBERTa} (\en{cross-encoder}) & -- & \checkmark & Ναι & Μία διέλευση ανά ζεύγος· τετραγωνική προσοχή στο μήκος εισόδου \\
        \en{Sentence-BERT} & -- & \checkmark & Ναι & Κωδικοποίηση ανά κόμβο $+$ $O(d)$ συνημίτονο ανά ερώτημα \\
        \en{LinkBERT} & \checkmark & \checkmark & Ναι (προ-εκπαίδευση) & Ισοδύναμο \en{BERT} ανά έγγραφο \\
        \en{GraphFormers} & \checkmark & \checkmark & Ναι, πλήρες γράφημα & Υψηλό (ΝΔΓ $+$ Μετασχηματιστής) \\
        \en{Patton} & \checkmark & \checkmark & Ναι (προ-εκπαίδευση) & Ισοδύναμο \en{BERT} $+$ όρος γείτονα \\
        \en{CascadeLP} (παρούσα εργασία) & \checkmark & \checkmark & Ναι, ανά επίπεδο & Μεταβλητό---κλιμακούμενο ανά ζεύγος \\
        \bottomrule
    \end{tabular}
    \caption{Σύγκριση οικογενειών μεθόδων πρόβλεψης ακμών σε γραφήματα με κειμενικά
    χαρακτηριστικά κόμβων.}
    \label{tab:method-family-comparison}
\end{table}


\section{Το Σύνολο Αξιολόγησης \en{Wikipedia DSAA 2023}}
\label{sec:dsaa-dataset}

Οι οικογένειες μεθόδων που παρουσιάστηκαν στις Ενότητες~\ref{sec:structural-methods}--\ref{sec:hybrid-architectures}---δομικοί ευρετικοί δείκτες, ενσωματώσεις γραφημάτων,
προ-εκπαιδευμένα γλωσσικά μοντέλα και υβριδικές αρχιτεκτονικές---χρειάζονται ένα
κοινό πεδίο αξιολόγησης για να συγκριθούν με νόημα.
Επιλέγουμε το σύνολο αξιολόγησης \en{Wikipedia DSAA 2023}~\cite{phan2023link,dsaa-2023-competition} για τον σκοπό αυτό, επειδή διαθέτει
δημόσιο πίνακα κατάταξης διαγωνισμού με υπαρκτές, δημοσιευμένες λύσεις από
πολλαπλές ομάδες---δηλαδή επιτρέπει απευθείας σύγκριση με τα αποτελέσματα των
παραπάνω οικογενειών μεθόδων, αντί για μεμονωμένες αναφορές βαθμολογιών χωρίς
κοινό πλαίσιο.
Το σύνολο αξιολόγησης εισήχθη ως κοινή εργασία στο 10ο Διεθνές
Συνέδριο \en{IEEE} για Επιστήμη Δεδομένων και Προηγμένα Αναλυτικά, μέσω
διαγωνισμού που φιλοξενήθηκε στο \en{Kaggle}.
Το σύνολο δεδομένων κατασκευάζεται από ένα στιγμιότυπο του γραφήματος
υπερσυνδέσμων της αγγλικής \en{Wikipedia} και σχεδιάζεται για τον έλεγχο μεθόδων
πρόβλεψης ακμών σε ένα γράφημα μεγάλης κλίμακας πλούσιο σε κείμενο.

Το σύνολο δεδομένων αποτελείται από τρία αρχεία.
Το \textbf{αρχείο εκπαίδευσης}
(\texttt{\en{train.csv}}) περιέχει \TrainPairs{} ζεύγη κόμβων με δυαδικές ετικέτες: \TrainPosPct\%
θετικά (υπάρχει ακμή) και \TrainNegPct\% αρνητικά (δεν υπάρχει ακμή).
Η ήπια ανισορροπία
κλάσεων δεν απαιτεί επιθετική επαναδειγματοληψία, αν και παρακινεί τη χρήση της
βαθμολογίας \en{Macro F1} ως μετρική αξιολόγησης αντί της ακρίβειας.
Το \textbf{αρχείο
κόμβων} (\texttt{\en{nodes.tsv}}) περιέχει μία γραμμή ανά κόμβο \en{Wikipedia} με
δύο στήλες: το αναγνωριστικό κόμβου και το ακατέργαστο κείμενο άρθρου.
Το κείμενο
μπορεί να περιλαμβάνει σήμανση \en{MediaWiki}, όπως περιεχόμενο \en{infobox}, αλλά
η υλοποίηση δεν το διαχωρίζει σε ανεξάρτητα πεδία.
Από τους \GraphNodesTotal{} κόμβους του συνόλου δεδομένων, \GraphNodes{}
εμφανίζονται στη δομική (γειτνιακή) συνιστώσα του γραφήματος· οι υπόλοιποι δεν έχουν
δομικούς γείτονες και αποτελούν τον πληθυσμό \en{cold-start}.
Το
\textbf{αρχείο ελέγχου} (\texttt{\en{test.csv}}) περιέχει \TestPairs{} ζεύγη κόμβων χωρίς
ετικέτες· η μετρική αξιολόγησης είναι η βαθμολογία \en{Macro F1} που υποβάλλεται
στη λίστα κατάταξης του διαγωνισμού.

Δύο συμμετοχές στο \en{DSAA 2023} αποτελούν τα κύρια μοντέλα αναφοράς της παρούσας
εργασίας.
Οι \textcite{phan2023link} (ομάδα \en{UIT-NLP}) διατυπώνουν την πρόβλεψη
ακμών ως \en{Natural Language Inference (NLI)}: συνενώνουν τον τίτλο και την περιγραφή
αμφότερων των κόμβων σε ένα ζεύγος \en{Premise-Hypothesis} και προσαρμόζουν ένα μοντέλο
\en{XLM-RoBERTa} για να προβλέψει συνεπαγωγή (θετική ακμή) ή ουδετερότητα-αντίφαση
(δεν υπάρχει ακμή).
Αυτή η προσέγγιση επιτυγχάνει βαθμολογία \en{Macro F1} 1.0 στο ιδιωτικό
σκέλος της λίστας κατάταξης, ισοβαθμώντας στην 3η θέση με τις δύο πρώτες.
Η
βαθμολογία αποδεικνύει επιτυχία στο συγκεκριμένο \en{leaderboard}, αλλά από μόνη
της δεν προσδιορίζει αν το αποφασιστικό σήμα προέρχεται από τη σχέση των δύο
άρθρων, από απλά χαρακτηριστικά ή από τεχνούργημα κατασκευής του συνόλου.
Οι \textcite{tran2023text} καταδεικνύουν
ότι ένα απλούστερο χαρακτηριστικό συχνότητας μερών του λόγου σε συνδυασμό με
ταξινομητή Τυχαίου Δάσους επιτυγχάνει βαθμολογία \en{Macro F1} 0.99999 (7η θέση),
παρέχοντας ένα σημαντικό στοιχείο ότι η βαθιά σημασιολογική κατανόηση δεν είναι αυστηρά
αναγκαία για αυτό το σύνολο αξιολόγησης.
Επιλέγονται ως κύρια μοντέλα αναφοράς επειδή
αντιπροσωπεύουν τα δύο άκρα του φάσματος κόστους (Μετασχηματιστής έναντι κλασικού
ταξινομητή) που εξετάζει το \en{CascadeLP}, όχι επειδή είναι οι μόνες δημόσια
τεκμηριωμένες συμμετοχές.

Ο διαγωνισμός απαιτούσε από τις κορυφαίες ομάδες να δημοσιεύσουν σύντομη αναφορά της
λύσης τους, και αρκετές ακόμα συμμετοχές είναι διαθέσιμες προς μελέτη· ο
Πίνακας~\ref{tab:dsaa-submissions} τις συνοψίζει.
Οι \textcite{nguyen2023mat} (ομάδα
\en{UIT Dark Cow}) προτείνουν το \en{Mutual Attention Transformer (MAT)}, έναν
μηχανισμό αμοιβαίας προσοχής μεταξύ των αναπαραστάσεων των δύο κόμβων, καταλαμβάνοντας
την 5η/4η θέση (δημόσιο/ιδιωτικό σκέλος).
Ο \textcite{yang2023achieving} συνδυάζει
τρεις ταξινομητές---\en{gradient boosting}, δασικό σύνολο δέντρων, και νευρωνικό
δίκτυο---με ψηφοφορία πλειοψηφίας, πετυχαίνοντας \en{Macro F1} άνω του 0.99 στην 6η
θέση.
Οι \textcite{giannoulidis2023enhanced} ακολουθούν διαφορετική στρατηγική:
αντί να επεξεργαστούν όλο το σύνολο δεδομένων, ανασυνθέτουν εν μέρει το γράφημα από τις
περιγραφές κόμβων, παράγουν διανύσματα ζεύγους βάσει μετρικών γραφήματος, και εφαρμόζουν
ένα κατώφλι ομοιότητας \en{TF-IDF}---μια προσέγγιση σχεδιασμένη ρητά για επεκτασιμότητα,
που επιτυγχάνει \en{Macro F1} 0.948, αισθητά χαμηλότερο από τις υπόλοιπες συμμετοχές
του Πίνακα~\ref{tab:dsaa-submissions}.
Οι \textcite{mata2023link} (ομάδα \en{Math-Land})
εκπαιδεύουν ενσωματώσεις τύπου \en{NLP} απευθείας εντός του νευρωνικού δικτύου
πρόβλεψης (χωρίς ξεχωριστό στάδιο προ-εκπαίδευσης ενσωμάτωσης), με απώλεια διασταυρούμενης
εντροπίας.
Οι \textcite{kansal2023predict}, τέλος, συνδυάζουν έναν αλγόριθμο αναζήτησης
κατά βάθος στο γράφημα με βαθμολογία κειμενικής ομοιότητας σε ένα ενιαίο μοντέλο
συνόλου, χωρίς να αναφέρουν συγκεκριμένη βαθμολογία ή θέση κατάταξης.

\begin{table}[htbp]
    \centering
    \small
    \begin{tabular}{@{}p{2.6cm}p{5.4cm}p{2.2cm}p{1.6cm}@{}}
        \toprule
        Ομάδα & Μέθοδος (περίληψη) & \en{Macro F1} & Θέση \\
        \midrule
        \en{UIT-NLP}~\cite{phan2023link} & \en{NLI} με \en{XLM-RoBERTa} σε ζεύγος τίτλου/περιγραφής & 1.0 (ιδιωτικό) & 3η (ισοβαθμία 1η--3η) \\
        \en{UIT Dark Cow}~\cite{nguyen2023mat} & \en{Mutual Attention Transformer} & -- & 5η/4η \\
        Yang~\cite{yang2023achieving} & Σύνολο \en{gradient boosting} $+$ δέντρα $+$ ΝΔ με ψηφοφορία & $>$0.99 & 6η \\
        Tran et al.~\cite{tran2023text} & Συχνότητα \en{POS} $+$ Τυχαίο Δάσος & 0.99999 & 7η \\
        Giannoulidis \& Mavroudopoulos~\cite{giannoulidis2023enhanced} & Μερική ανασύνθεση γραφήματος $+$ κατώφλι ομοιότητας \en{TF-IDF} & 0.948 & -- \\
        \en{Math-Land}~\cite{mata2023link} & Ενσωματώσεις τύπου \en{NLP} εκπαιδευμένες \en{end-to-end} & -- & -- \\
        Kansal \& Mehta~\cite{kansal2023predict} & Αναζήτηση κατά βάθος $+$ ομοιότητα κειμένου (σύνολο) & -- & -- \\
        \bottomrule
    \end{tabular}
    \caption{Δημόσια τεκμηριωμένες συμμετοχές στο \en{DSAA 2023 Wikipedia Link
    Prediction}. Παύλα (\en{--}) σημαίνει ότι η αντίστοιχη αναφορά δεν δηλώνει ρητά
    την τιμή.}
    \label{tab:dsaa-submissions}
\end{table}

Το μοτίβο του Πίνακα~\ref{tab:dsaa-submissions} ενισχύει την παρατήρηση που παρακίνησε
αυτή την εργασία: όλες οι συμμετοχές που δηλώνουν αριθμητική βαθμολογία ξεπερνούν το
0.94, και οι περισσότερες ξεπερνούν το 0.99, ανεξάρτητα από το εάν η μέθοδος αξιοποιεί
βαθύ Μετασχηματιστή (\en{UIT-NLP}), κλασικό σύνολο ταξινομητών (\en{Yang}), ή απλά
χαρακτηριστικά μερών του λόγου (\en{Tran et al.}).
Αυτή η σύγκλιση ανεξαρτήτως
μεθοδολογικής πολυπλοκότητας είναι ακριβώς το είδος του σήματος που παρακινεί τον
ρητό έλεγχο διαρροής και διαχωρισιμότητας της Ενότητας~\ref{sec:leakage-background},
αντί της απευθείας αποδοχής της βαθμολογίας ως απόδειξη μεθοδολογικής ανωτερότητας.


\section{Ακεραιότητα Αξιολόγησης και Διαρροή Δεδομένων}
\label{sec:leakage-background}

Σχεδόν τέλειες βαθμολογίες σε ένα σύνολο αξιολόγησης πρόβλεψης ακμών μπορούν να
προκύψουν από τρεις διακριτές αιτίες, οι οποίες συχνά συγχέονται στη βιβλιογραφία:
γνήσια ισχυρή μεθοδολογία, διαρροή δεδομένων μεταξύ εκπαίδευσης και ελέγχου, ή
εγγενής διαχωρισιμότητα του συνόλου δεδομένων (δηλαδή, ένα μεγάλο κλάσμα ζευγών
είναι τετριμμένα ταξινομήσιμο ανεξάρτητα από τη μέθοδο).
Η διάκριση μεταξύ αυτών
των τριών αιτιών απαιτεί ρητό έλεγχο, όχι μόνο αναφορά συνολικής μετρικής — μια
ανησυχία που έχει τεκμηριωθεί και σε ευρέως χρησιμοποιούμενα σημεία αναφοράς όπως
το \en{OGB}, όπου προβλήματα στη σχεδίαση της διαίρεσης δεδομένων και στη
δειγματοληψία αρνητικών παραδειγμάτων έχουν οδηγήσει σε παραπλανητικές συγκρίσεις
μεταξύ μεθόδων~\cite{li2023evaluating}.

Δύο μορφές διαρροής είναι ιδιαίτερα συναφείς για σύνολα δεδομένων πρόβλεψης ακμών
που κατασκευάζονται με τυχαία δειγματοληψία θετικών και αρνητικών ζευγών από ένα
ενιαίο γράφημα.
Η πρώτη είναι η \textit{επικάλυψη ζευγών}: εάν το ίδιο ζεύγος (ή το
αντεστραμμένο του, σε μη κατευθυνόμενο γράφημα) εμφανίζεται τόσο στο σύνολο
εκπαίδευσης όσο και στο σύνολο ελέγχου, ένα μοντέλο θα μπορούσε να απομνημονεύσει
την ετικέτα αντί να τη συμπεράνει.
Η δεύτερη είναι η \textit{ενδο-συνόλου
διπλοεγγραφή}: εάν το ίδιο ζεύγος εμφανίζεται πάνω από μία φορά εντός του συνόλου
εκπαίδευσης, μια τυχαία διαίρεση εκπαίδευσης-επαλήθευσης βάσει γραμμής (αντί βάσει
ταυτότητας ζεύγους) μπορεί να τοποθετήσει το ίδιο φυσικό ζεύγος και στις δύο
πλευρές της διαίρεσης, διογκώνοντας τεχνητά τη βαθμολογία επαλήθευσης.

Η παρούσα εργασία υιοθετεί μια ρητή πρωτόκολλη ελέγχου διαρροής (Κεφάλαιο~\ref{ch:methodology})
για το σύνολο \en{DSAA 2023}, ελέγχοντας αμφότερες τις
μορφές διαρροής πριν από την ερμηνεία οποιασδήποτε αναφερόμενης βαθμολογίας.
Αυτό αντιμετωπίζει ένα κενό στη βιβλιογραφία για αυτό το σύνολο αξιολόγησης: ούτε
οι \textcite{phan2023link} ούτε οι \textcite{tran2023text} αναφέρουν
έλεγχο διαρροής μεταξύ των αρχείων \texttt{\en{train.csv}} και
\texttt{\en{test.csv}}.


\section{Κενά στην Υπάρχουσα Εργασία}
\label{sec:gaps}

Η υπάρχουσα βιβλιογραφία για αυτό το σύνολο αξιολόγησης και για την πρόβλεψη ακμών
σε γραφήματα με κειμενικά χαρακτηριστικά κόμβων γενικότερα αφήνει αρκετά σημαντικά
κενά που αντιμετωπίζει η παρούσα εργασία.

\textbf{Διαφάνεια κόστους.} Τα προηγούμενα κορυφαία συστήματα αναφέρουν ακρίβεια
αλλά όχι κόστος συμπερασμού.
Η παραγωγή ενσωματώσεων \en{Sentence-Transformer} είναι ακριβότερη από την
αναζήτηση προϋπολογισμένων δομικών χαρακτηριστικών, αλλά το πραγματικό κόστος
εξαρτάται από το αν οι κόμβοι κωδικοποιούνται εκ των προτέρων ή κατ' απαίτηση.
Το \en{CascadeLP} σχεδιάζεται ώστε να
καθιστά αυτόν τον συμβιβασμό ρητό, μετρώντας το κλάσμα ζευγών που επιλύεται σε κάθε
επίπεδο και τον ρυθμό επεξεργασίας της αλυσίδας πάνω σε προϋπολογισμένα
χαρακτηριστικά.
Η πραγματική εξοικονόμηση κωδικοποίησης παραμένει υπόθεση για
υλοποίηση κατ' απαίτηση και δεν εξισώνεται με το ποσοστό δρομολόγησης.

\textbf{Ανάλυση \en{cold-start}.} Τόσο οι Phan et al.\ όσο και οι Tran
et al.\ αναφέρουν συνολικές βαθμολογίες χωρίς αποδιαχωρισμό της απόδοσης σε ζεύγη
\en{cold-start} (κόμβοι χωρίς δομικούς γείτονες).
Στο σύνολο δεδομένων
\en{DSAA 2023}, οι κόμβοι \en{cold-start} αποτελούν σημαντικό κλάσμα του
πληθυσμού κόμβων (Κεφάλαιο~\ref{ch:methodology}), και ένα σύστημα που τα
αντιμετωπίζει ανεπαρκώς θα υποβαθμιστεί καθώς εξελίσσεται το γράφημα.
Το Επίπεδο~2 του
\en{CascadeLP} σχεδιάζεται ειδικά για την περίπτωση \en{cold-start} και
παρακάμπτει ρητά το δομικό επίπεδο.
Η πρόσθετη ανάλυση αναφέρει χωριστά το
ευρύτερο υποσύνολο μηδενικού \en{CN}, χωρίς να το ταυτίζει με τον λειτουργικό
ορισμό της δρομολόγησης.

\textbf{Χαρακτηρολόγηση συνόλου δεδομένων.} Οι σχεδόν τέλειες βαθμολογίες που
αναφέρονται για αυτό το σύνολο αξιολόγησης εγείρουν το ερώτημα εάν το σύνολο
αξιολόγησης είναι κορεσμένο ή εάν περιέχει μεγάλο κλάσμα τετριμμένα ταξινομήσιμων
ζευγών.
Η παρούσα εργασία παρέχει συστηματική χαρακτηρολόγηση, εντοπίζοντας \SelfLoopsTrain{}
αυτοβρόχους στο σύνολο εκπαίδευσης (\SelfLoopsTrainPos{} θετικοί, τετριμμένα προβλέψιμοι χωρίς
οποιοδήποτε μοντέλο, και \SelfLoopsTrainNeg{} αρνητικός --- μια ανωμαλία ετικέτας), αναλύοντας
την κατανομή πλήθους κοινών γειτόνων και ποσοτικοποιώντας το υποσύνολο
μηδενικού \en{CN}.
Αυτή η χαρακτηρολόγηση πλαισιώνει τις αναφερόμενες βαθμολογίες και
είναι απαραίτητη για την κατανόηση της αφετηρίας της μελλοντικής προόδου σε αυτό το
σύνολο αξιολόγησης.

\textbf{Έλεγχος διαδικασίας δειγματοληψίας.} Οι προηγούμενες αναφορές δεν
εξετάζουν αν η ταυτότητα ενός άκρου προβλέπει σχεδόν μόνη της την ετικέτα.
Η
παρούσα εργασία ελέγχει ρητά την κατανομή των αρνητικών ως προς τα δύο άκρα και
συγκρίνει τα αποτελέσματα του \en{DSAA 2023} με νέο σύνολο στο οποίο οι αρνητικές
ακμές δειγματοληπτούνται από τον ίδιο πληθυσμό κόμβων.

\textbf{Αποτελεσματικότητα υβριδικών μεθόδων.} Υφιστάμενες υβριδικές αρχιτεκτονικές
όπως τα \en{GraphFormers} και το \en{Patton} είναι δαπανηρές στην εκπαίδευση και
ανάπτυξη.
Το \en{CascadeLP} υιοθετεί διαφορετική προσέγγιση: αντί να κατασκευάζει ένα
μονολιθικό μοντέλο που συνδυάζει όλα τα σήματα κατά την εκπαίδευση, διατάσσει απλά
μοντέλα σε σειρά και κλιμακώνει σε ακριβότερα μόνο όταν είναι αναγκαίο.
Αυτός ο
σχεδιασμός είναι ερμηνεύσιμος, αρθρωτός και εύκολα επεκτάσιμος όταν νέες οικογένειες
χαρακτηριστικών καθίστανται διαθέσιμες.
